\documentclass[aspectratio=169, 10pt]{beamer}

\usepackage[utf8]{inputenc}
\usepackage[T1]{fontenc}
\usepackage[spanish]{babel}
\usepackage{amsmath, amssymb}
\usepackage{microtype}
\usepackage{booktabs}
\usepackage{colortbl}
\usepackage{siunitx}
\usepackage[american, RPvoltages]{circuitikz}

\sisetup{output-decimal-marker = {,}, per-mode = symbol}

\definecolor{ucbblue}{RGB}{0, 51, 102}
\usetheme{Boadilla}
\usecolortheme{seahorse}
\setbeamertemplate{navigation symbols}{}
\setbeamertemplate{itemize items}[circle]
\setbeamercolor{title}{bg=ucbblue, fg=white}
\setbeamercolor{frametitle}{bg=ucbblue, fg=white}

\title[Análisis Nodal Canónico]{Análisis Nodal Canónico ($\mathbf{G}\mathbf{V}=\mathbf{I}$)}
\subtitle{Supernodos y Modelado Modular en LTspice}
\author[M.Sc. Ing. B. Quiroga]{M.Sc. Ing. Bernardo Quiroga Turdera}
\institute[UCB Tarija]{IMT-121 Circuitos Electrónicos I | Semestre II/2026}
\date{Sesión 09}

\begin{document}

\begin{frame}
    \titlepage
\end{frame}

\begin{frame}{Resultados de Aprendizaje (Criterios ABET 1 y 6)}
    \begin{block}{Objetivos de la Sesión}
        \begin{itemize}
            \item \textbf{Construir} por inspección directa la matriz de conductancias $\mathbf{G}$.
            \item \textbf{Formular y resolver} supernodos con fuentes de tensión flotantes.
            \item \textbf{Estructurar} esquemáticos modulares jerárquicos y parametrizados en LTspice.
        \end{itemize}
    \end{block}
\end{frame}

\begin{frame}{El Principio de Dualidad (De Mallas a Nodos)}
    \textbf{¿Por qué cambiar de enfoque?}
    \begin{itemize}
        \item \textbf{Mallas:} Analiza corrientes de lazo usando resistencias ($\mathbf{R}\mathbf{I}=\mathbf{V}$).
        \item \textbf{Nodos:} Analiza potenciales respecto a Tierra usando conductancias ($\mathbf{G}\mathbf{V}=\mathbf{I}$).
    \end{itemize}
    \textit{Conductancia ($G$):} $G = \frac{1}{R}$ expresada en Siemens (\unit{\siemens}).

    \vspace{0.4cm}
    \begin{alertblock}{Reglas Mnemotécnicas de Inspección Directa ($\mathbf{G}\mathbf{V}=\mathbf{I}$)}
        \begin{itemize}
            \item \textbf{$G_{kk}$ (Diagonal):} Suma de conductancias adheridas al nodo $k$ ($>0$).
            \item \textbf{$G_{kj}$ (Fuera Diagonal):} Negativo de la conductancia compartida entre $k$ y $j$ ($\le 0$).
            \item \textbf{$I_{sk}$ (Vector de Fuentes):} Suma algebraica de corrientes independientes que \textit{entran} al nodo $k$.
        \end{itemize}
    \end{alertblock}
\end{frame}

\begin{frame}{El Tratamiento de Supernodos}
    Ocurre cuando una fuente de tensión ideal une dos nodos \textbf{no referenciados a Tierra} (flotantes). La inspección directa colapsa porque $R \to 0 \implies G \to \infty$.

    \begin{columns}
        \begin{column}{0.45\textwidth}
            \begin{center}
                \begin{circuitikz}[scale=0.8, transform shape]
                    \draw (0,0) node[left]{Nodo $a$} to[short, o-] (1.5,0)
                          to[V, v=$V_s$] (4.5,0)
                          to[short, -o] (6,0) node[right]{Nodo $b$};
                    \draw[dashed, red, thick, rounded corners] (0.5, -1) rectangle (5.5, 1);
                    \node[red, above] at (3, 1) {\textbf{SUPERNODO}};
                \end{circuitikz}
            \end{center}
        \end{column}
        \begin{column}{0.55\textwidth}
            \textbf{Protocolo Obligatorio (Dos Ecuaciones):}
            \begin{enumerate}
                \item \textbf{Ecuación LCK (Conservación):} Sumar todas las corrientes que escapan de la frontera del supernodo ($= 0$).
                \item \textbf{Ecuación LVK (Restricción):} Relación de la fuente de tensión:
                $$V_{\text{positivo}} - V_{\text{negativo}} = V_{\text{fuente}}$$
            \end{enumerate}
        \end{column}
    \end{columns}
\end{frame}

\begin{frame}{Resumen de Resultados del Ejemplo Integrador}
    \textbf{Circuito Base:} $I_{s1} = \qty{3}{\ampere}$, $I_{s2} = \qty{1}{\ampere}$, $V_s = \qty{6}{\volt}$, $R_1=R_4=\qty{2}{\ohm}$, $R_2=R_3=\qty{4}{\ohm}$.
    
    \vspace{0.4cm}
    \textbf{Convergencia Tripartita de Calidad (ABET 6):}
    \begin{table}[]
        \centering
        \renewcommand{\arraystretch}{1.4}
        \begin{tabular}{@{}l c c c c@{}}
            \toprule
            \rowcolor{ucbblue!10}
            \textbf{Variable Nodal} & \textbf{Teoría Analítica} & \textbf{Script Python} & \textbf{LTspice \texttt{.op}} & \textbf{Error Relativo ($\varepsilon_r$)} \\
            \midrule
            $V_1$ & \qty{5.8182}{\volt} & \qty{5.8182}{\volt} & \qty{5.8182}{\volt} & $0.00\%$ \\
            $V_2$ & \qty{5.4545}{\volt} & \qty{5.4545}{\volt} & \qty{5.4545}{\volt} & $0.00\%$ \\
            $V_3$ & \qty{-0.5455}{\volt} & \qty{-0.5455}{\volt} & \qty{-0.5455}{\volt} & $0.00\%$ \\
            \bottomrule
        \end{tabular}
    \end{table}
\end{frame}

\end{document}
